
\documentclass{article}
\usepackage{colortbl}
\usepackage{makecell}
\usepackage{multirow}
\usepackage{supertabular}

\begin{document}

\newcounter{utterance}

\centering \large Interaction Transcript for game `eqbench', experiment `eq\_bench', episode 135 with refineonly{-}Qwen3.5{-}4B.
\vspace{24pt}

{ \footnotesize  \setcounter{utterance}{1}
\setlength{\tabcolsep}{0pt}
\begin{supertabular}{c@{$\;$}|p{.15\linewidth}@{}p{.15\linewidth}p{.15\linewidth}p{.15\linewidth}p{.15\linewidth}p{.15\linewidth}}
    \# & \multicolumn{2}{c}{Player} && \multicolumn{2}{c}{Game Master} \\
    \hline

    \theutterance \stepcounter{utterance}  
    & & & \multicolumn{4}{p{0.6\linewidth}}{
        \cellcolor[rgb]{0.9,0.9,0.9}{
            \makecell[{{p{\linewidth}}}]{
                \texttt{\tiny{[P1$\langle$GM]}}
                \\ 
\texttt{Your task is to predict the likely emotional responses of a character in this dialogue:} \\
\\ 
\texttt{Jacques: Alessandra, you've been under{-}seasoning the scallops all night. } \\
\texttt{Alessandra: I've been following the recipe, chef.} \\
\texttt{Jacques: (Laughs) The recipe? That's your problem, Alessandra. You're a slave to the recipe. You need to feel the food, taste it, not just blindly follow instructions.} \\
\texttt{Alessandra: I understand that, chef.} \\
\texttt{Jacques: Do you? Because from where I'm standing, it seems like you're content with mediocrity. } \\
\texttt{Alessandra: That's not fair, Jacques.} \\
\texttt{Jacques: Fair? This isn't about fairness, Alessandra. This is about excellence. } \\
\texttt{Alessandra: I'm doing my best.} \\
\texttt{Jacques: Your best isn't good enough. Not for this kitchen. Not for me.} \\
\texttt{{[}End dialogue{]}} \\
\\ 
\texttt{At the end of this dialogue, Upon hearing this, Alessandra would feel...} \\
\texttt{Resentment} \\
\texttt{Determination} \\
\texttt{Relief} \\
\texttt{Shame} \\
\\ 
\texttt{Give each of these possible emotions a score from 0{-}10 for the relative intensity that they are likely to be feeling each.} \\
\\ 
\texttt{You must output in the following format, including headings (of course, you should give your own scores), with no additional commentary:} \\
\\ 
\texttt{Resentment: <score>} \\
\texttt{Determination: <score>} \\
\texttt{Relief: <score>} \\
\texttt{Shame: <score>} \\
\\ 
\\ 
\texttt{{[}End of answer{]}} \\
\\ 
\texttt{Remember: zero is a valid score, meaning they are likely not feeling that emotion. You must score at least one emotion > 0.} \\
\\ 
\texttt{Your answer:} \\
            }
        }
    }
     \\ \\

    \theutterance \stepcounter{utterance}  
    & & \multicolumn{4}{p{0.6\linewidth}}{
        \cellcolor[rgb]{0.95,0.95,0.95}{
            \makecell[{{p{\linewidth}}}]{
                \texttt{\tiny{[P1$\leftrightarrow$P]}}
                \texttt{{[}DRAFT{]}} \\
\texttt{Resentment: 4} \\
\texttt{Determination: 8} \\
\texttt{Relief: 0} \\
\texttt{Shame: 6} \\
\\ 
\texttt{**Reasoning:**} \\
\texttt{*   **Determination (8):** Alessandra has consistently defended her approach ("I've been following the recipe," "I understand that," "I'm doing my best"). Despite Jacques's harsh criticism and dismissal of her efforts, she maintains her stance rather than giving up, indicating a strong resolve to prove herself or continue her current path.} \\
\texttt{*   **Shame (6):** Jacques explicitly attacks her character and work ethic, calling her a "slave to the recipe" and accusing her of being "content with mediocrity." This public humiliation and the implication that she is failing to meet high standards would naturally induce feelings of shame or inadequacy.} \\
\texttt{*   **Resentment (4):** While she is being criticized, she hasn't yet lashed out or shown anger; she is still trying to explain herself ("That's not fair," "I'm doing my best"). Resentment is a likely next step but is currently lower than determination and shame.} \\
\texttt{*   **Relief (0):** There is no indication of tension release or a problem being solved; the situation is escalating in conflict.} \\
            }
        }
    }
    & \\ \\

    \theutterance \stepcounter{utterance}  
    & \multicolumn{4}{p{0.6\linewidth}}{
        \cellcolor[rgb]{0.561,0.737,0.561}{
            \makecell[{{p{\linewidth}}}]{
                \texttt{\tiny{[P1$\rangle$GM]}}
                \texttt{Resentment: 5} \\
\texttt{Determination: 7} \\
\texttt{Relief: 0} \\
\texttt{Shame: 7} \\
\\ 
\texttt{**Reasoning:**} \\
\texttt{*   **Shame (7):** Jacques's attack on Alessandra's character ("slave to the recipe") and her work ethic ("content with mediocrity") is a direct blow to her self{-}esteem. In a high{-}stakes environment like a kitchen, being publicly called out for mediocrity by a superior often triggers a strong sense of shame or inadequacy, making this the most dominant immediate emotion.} \\
\texttt{*   **Determination (7):** While she defends herself, the intensity of the criticism ("Your best isn't good enough. Not for this kitchen. Not for me") suggests she is feeling the weight of the failure. Her determination is high, but the emotional impact of the insult likely overshadows it slightly in the immediate moment, though it remains very strong.} \\
\texttt{*   **Resentment (5):** The repeated dismissal and the phrase "That's not fair" indicate the beginning of a buildup of anger and resentment. However, she hasn't yet crossed the threshold into open hostility or defiance, so it is currently lower than shame and determination.} \\
\texttt{*   **Relief (0):** The conflict is escalating, not resolving, so relief remains at zero.} \\
            }
        }
    }
    & & \\ \\

    \theutterance \stepcounter{utterance}  
    & & & \multicolumn{2}{p{0.3\linewidth}}{
        \cellcolor[rgb]{0.9,0.9,0.9}{
            \makecell[{{p{\linewidth}}}]{
                \texttt{\tiny{[GM$|$GM]}}
                \texttt{Resentment: 5} \\
\texttt{Determination: 7} \\
\texttt{Relief: 0} \\
\texttt{Shame: 7} \\
            }
        }
    }
    & & \\ \\

    \theutterance \stepcounter{utterance}  
    & & & \multicolumn{2}{p{0.3\linewidth}}{
        \cellcolor[rgb]{0.9,0.9,0.9}{
            \makecell[{{p{\linewidth}}}]{
                \texttt{\tiny{[GM$|$GM]}}
                \texttt{Resentment: 8} \\
\texttt{Determination: 6} \\
\texttt{Relief: 0} \\
\texttt{Shame: 7} \\
            }
        }
    }
    & & \\ \\

    \theutterance \stepcounter{utterance}  
    & & & \multicolumn{2}{p{0.3\linewidth}}{
        \cellcolor[rgb]{0.9,0.9,0.9}{
            \makecell[{{p{\linewidth}}}]{
                \texttt{\tiny{[GM$|$GM]}}
                \texttt{game\_result = LOSE} \\
            }
        }
    }
    & & \\ \\

\end{supertabular}
}

\end{document}
